\documentclass[12pt, a4paper]{article}
\usepackage[UTF8, zihao=-4]{ctex}
\usepackage[top=2.5cm, bottom=2.5cm, left=2.8cm, right=2.8cm]{geometry}
\usepackage{fontspec}
\setmainfont{Times New Roman}
\usepackage{setspace}
\setstretch{1.5}
\usepackage{booktabs}
\usepackage{array}
\usepackage{makecell}
\usepackage{multirow}
\usepackage{adjustbox}
\usepackage{amsmath, amssymb}
\usepackage{fancyhdr}
\usepackage{float}
\setlength{\headheight}{15pt}
\pagestyle{fancy}
\fancyhf{}
\fancyhead[C]{\small 实验数据记录单}
\fancyfoot[C]{\small\thepage}

\begin{document}

% ===== 封面信息 =====
\begin{center}
  {\heiti\zihao{3} 薄透镜焦距三种测量方法 \\ 实验数据记录单} \\[1cm]
  \begin{tabular}{llll}
    姓\quad 名：& \underline{\hspace{4cm}} &
    学\quad 号：& \underline{\hspace{4cm}} \\[0.5cm]
    班\quad 级：& \underline{\hspace{4cm}} &
    实验日期：  & \underline{\hspace{4cm}} \\[0.5cm]
    指导教师：  & \underline{\hspace{4cm}} &
    教师签名：  & \underline{\hspace{4cm}} \\
  \end{tabular}
\end{center}
\vspace{0.8cm}

% ===== 表1：自准直法 =====
\subsection*{表 1 \quad 自准直法测量薄透镜焦距记录表}
{\footnotesize 仪器：F形LED光源、待测透镜（$\Phi50$，标称焦距 $f=200\,\text{mm}$）、平面反射镜；测量公式：$f = a_1 - a_2$}

\begin{table}[H]
  \centering
  \begin{adjustbox}{max width=\textwidth}
    \begin{tabular}{>{\centering\arraybackslash}m{2cm}
                    >{\centering\arraybackslash}m{4cm}
                    >{\centering\arraybackslash}m{4cm}
                    >{\centering\arraybackslash}m{4cm}}
      \toprule
      测量次数 & 目标板位置 $a_1$ / mm & 被测透镜位置 $a_2$ / mm & 透镜焦距 $f$ / mm \\
      \midrule
      1 & & & \\[0.6cm]
      2 & & & \\[0.6cm]
      3 & & & \\[0.6cm]
      \midrule
      \multicolumn{3}{c}{焦距平均值 $\bar{f}$ / mm} & \\[0.6cm]
      \bottomrule
    \end{tabular}
  \end{adjustbox}
\end{table}

\vspace{0.5cm}

% ===== 表2：二次成像法 =====
\subsection*{表 2 \quad 二次成像法（贝塞尔法）测量薄透镜焦距记录表}
{\footnotesize 仪器：F形LED光源、待测透镜（$\Phi50$，标称焦距 $f=150\,\text{mm}$）、白屏；物屏距 $l \geq 600\,\text{mm}$；\\
测量公式：$d = a_2 - a_1$，\quad $f' = \dfrac{l^2 - d^2}{4l}$}

\begin{table}[H]
  \centering
  \begin{adjustbox}{max width=\textwidth}
    \begin{tabular}{>{\centering\arraybackslash}m{1.8cm}
                    >{\centering\arraybackslash}m{3.5cm}
                    >{\centering\arraybackslash}m{3.5cm}
                    >{\centering\arraybackslash}m{2.5cm}
                    >{\centering\arraybackslash}m{2.5cm}
                    >{\centering\arraybackslash}m{2.8cm}}
      \toprule
      测量次数 &
      \makecell{第一次清晰\\成像位置\\$a_1$ / mm} &
      \makecell{第二次清晰\\成像位置\\$a_2$ / mm} &
      \makecell{两次位置差\\$d = a_2 - a_1$\\/ mm} &
      \makecell{物屏距\\$l$ / mm} &
      \makecell{透镜焦距\\$f'$ / mm} \\
      \midrule
      1 & & & & & \\[0.8cm]
      2 & & & & & \\[0.8cm]
      3 & & & & & \\[0.8cm]
      \midrule
      \multicolumn{5}{c}{焦距平均值 $\bar{f}'$ / mm} & \\[0.6cm]
      \bottomrule
    \end{tabular}
  \end{adjustbox}
\end{table}

\vspace{0.5cm}

% ===== 表3：焦距仪法 =====
\subsection*{表 3 \quad 平行光管法（焦距仪）测量薄透镜焦距记录表}
{\footnotesize 仪器：LED光源、平行光管（$f_o'=400\,\text{mm}$，玻罗板最大线间距 $y=4\,\text{mm}$）、CMOS相机（像元 $2.9\,\mu\text{m}$）；\\
测量公式：$f_x' = \dfrac{y'}{y} \times f_o'$，软件直接输出焦距值（单位 mm）}

\begin{table}[H]
  \centering
  \begin{adjustbox}{max width=\textwidth}
    \begin{tabular}{>{\centering\arraybackslash}m{4.5cm}
                    >{\centering\arraybackslash}m{2.5cm}
                    >{\centering\arraybackslash}m{2.5cm}
                    >{\centering\arraybackslash}m{2.5cm}
                    >{\centering\arraybackslash}m{3cm}}
      \toprule
      透镜 & 第1次 / mm & 第2次 / mm & 第3次 / mm & 平均值 / mm \\
      \midrule
      透镜1（标称 $200\,\text{mm}$） & & & & \\[0.8cm]
      透镜2（标称 $150\,\text{mm}$） & & & & \\[0.8cm]
      \bottomrule
    \end{tabular}
  \end{adjustbox}
\end{table}

\vspace{1cm}

% ===== 计算与结论区域 =====
\subsection*{计算与结论}

\noindent 三种方法测量焦距（$f = 200\,\text{mm}$ 透镜）汇总：

\vspace{0.3cm}
\begin{table}[H]
  \centering
  \begin{tabular}{>{\centering\arraybackslash}m{5cm}
                  >{\centering\arraybackslash}m{4cm}
                  >{\centering\arraybackslash}m{4cm}}
    \toprule
    方法 & 平均焦距 $\bar{f}$ / mm & 与标称值偏差 / \% \\
    \midrule
    自准直法 & & \\[0.6cm]
    二次成像法（贝塞尔法）& & \\[0.6cm]
    平行光管法 & & \\[0.6cm]
    \bottomrule
  \end{tabular}
\end{table}

\vspace{2cm}

\noindent\textbf{实验结论：}

\vspace{4cm}

\end{document}
